\section{What recovers: content or format}
\label{sec:mechanism}

Does recovery use the \emph{content} of the echoed registry, or only the
\emph{format} of a structured rejection? The condition ladder answers this
directly, without the distractor matrix we are underpowered to read.

\paragraph{The ablation already points away from content.}
On Qwen3-8B the structured error $\mathrm{C}_{0.7}$, a
\texttt{tool\_not\_found} envelope with \emph{no registry attached},
already drives fabrications to zero, matching full RVR
(Section~\ref{sec:results:rvr}). If the model needed to \emph{read} the
registry to recover, removing it should cost something; it costs
nothing.

\paragraph{The decoy arm closes the question.}
$\mathrm{C}_{0.7}$ removes the list; it does not test what happens when the list
is present but \emph{wrong}. A registry-as-cue account predicts a decoy list
(correctly formatted, count-matched, but populated with non-existent tools)
should fail to recover, with no valid name to copy. A format-only account
predicts it recovers anyway, since the constraint ``your last call was rejected;
choose from this set'' need not be truthful to suppress the out-of-registry
call. The $\mathrm{C}_{0.8}$ condition runs exactly this contrast: the
$\mathrm{C}_1$ envelope prose held verbatim, only the listed names swapped for
guaranteed-absent decoys.

The decoy list recovers at $0/410$ (Clopper--Pearson 95\% upper bound
$\leq\!0.73\%$) across all four probe arms (near\_name $0/97$, synonym
$0/106$, matched\_random $0/107$, unrelated $0/100$), indistinguishable
from the real list ($\mathrm{C}_1$, $0/258$ on 8B) and from the no-list
envelope ($\mathrm{C}_{0.7}$, $0/448$). A wrong list recovers as well as
the right one: registry \emph{content} is decorative for recovery. The
model is responding to the message-level rejection itself, not being
cued by valid names; the truthfulness of the offered set is not
load-bearing. This is positive evidence against the registry-as-cue
account, which would need a valid name to point at and the decoy offers
none.

\paragraph{What this does not claim.}
Format-sufficiency is a claim about \emph{recovery}, not task success. A decoy
list suppresses the fabricated call but offers no valid target, so where the real
tool was removed the model must still decline or redirect; the value of the
\emph{true} registry shows up in task-completion rate, not in TFR, which is the
quantity the rejection signal governs. Nor does format-sufficiency hold at every
scale: it is established only on the tiers where C0 fabrications exist to be
suppressed (1.7--14B); the 0.6B and 32B tiers carry no signal to recover.

\paragraph{The scale curve, descriptively.}
The dominant distractor type drifts with scale
(Section~\ref{sec:results:scaling}), but a Friedman exact permutation test is
underpowered ($\chi^2_{(3)}\!=\!3.0$, $p\!=\!0.46$; $n\!=\!4$ sizes), so we read
the drift as phenomenology, not inference. The content-vs-format question rests
not on this matrix but on the powered $\mathrm{C}_{0.7}$ and $\mathrm{C}_{0.8}$
ablations.
